\subsection{Evaluación Financiera}

\subsubsection{Márgenes de contribución}

Los ingresos se dividen en ventas de licencias e implementaciones. A continuación se calculan los márgenes de contribución de cada línea de negocio.

\paragraph{Margen de contribución por Implementación}

El margen bruto se calcula restando los costos de implementación al precio de venta:

\begin{quote}
Precio de venta por implementación: USD 10,000\\
Costos totales por implementación: USD 3,900\\
Margen bruto absoluto por implementación: USD 6,100
\end{quote}
	
\paragraph{Margen de contribución por licencia}

\begin{quote}
Precio de venta por licencia: USD 12,000\\
Costos variables asociados a gastos generales: USD 120\\
Costos variables asociados a marketing: USD 600\\
Margen absoluto por licencia anual: USD 11,280
\end{quote}

\subsubsection{Tasas para la valuación del proyecto}

En diciembre de 2024, las tasas de interés para préstamos a PyMEs en Argentina varían según la entidad financiera y las condiciones específicas del crédito. A continuación, se detallan algunas de las tasas ofrecidas por diferentes instituciones:

\begin{quote}
Banco Nación: Ofrece una Tasa Nominal Anual (TNA) del 34\% para créditos a 90 días, y 35\% para plazos superiores.\\
Banco Provincia: Presenta tasas del 40\% TNA para capital de trabajo a 12 meses.\\
Bancos Privados: Las tasas para préstamos de corto plazo (180 días) rondan el 49.5\% TNA, mientras que para plazos de 12 meses ascienden al 51\% TNA.
\end{quote}

Considerando estas opciones, una tasa promedio razonable para el análisis financiero sería del 40\% TNA. 

Esto significa que la TEA es aproximadamente del 49.27\% para intereses compuestos mensualmente.

\paragraph{Cálculo de la tasa de descuento con el método CAPM}

Según el método CAPM (Capital Asset Pricing Model), el costo del capital (Ke) se calcula usando la siguiente fórmula:

\[Ke = Rf + \beta (Rm - Rf) + \text{Riesgo País}\]

\textbf{Selección de la Tasa Libre de Riesgo}

La tasa libre de riesgo (Rf) se toma del rendimiento a 10 años de los bonos del tesoro de los Estados Unidos de Norteamérica, hoy en torno al 4,6\%.

\[Rf = 4,6\%\]

\textbf{Estimación del factor Beta}

Para este estudio, se propone un valor de Beta de 1.5, fundamentado en las características del sector, la naturaleza del producto y las condiciones del mercado regional.

El mercado tecnológico, particularmente en el ámbito de software empresarial y soluciones de inteligencia artificial, se caracteriza por una volatilidad superior al promedio del mercado general. Empresas maduras como Microsoft y Salesforce operan con valores de Beta entre 1 y 1.3, reflejando un riesgo moderado. Las startups tecnológicas que ingresan en mercados competitivos y en rápida evolución presentan Betas superiores, generalmente en el rango de 1.3 a 2.0, debido a su mayor dependencia del crecimiento y la incertidumbre inicial.

El middleware de hiperpersonalización que constituye el producto central de este proyecto está diseñado para integrarse con sistemas empresariales como CRMs, ERPs y bases de datos. Esto genera una menor volatilidad relativa en comparación con startups más disruptivas que operan en mercados especulativos como el de las criptomonedas. Sin embargo, el hecho de que sea un producto nuevo, con una curva de adopción inicial y competencia directa de actores consolidados (e.g., Salesforce, Amazon) aumenta el nivel de riesgo, justificando un Beta superior a 1.

La operación se desarrolla en un mercado emergente, América Latina, donde factores económicos y políticos incrementan la incertidumbre. Esto se traduce en un ajuste adicional al valor de Beta.

Según Damodaran (2024), las empresas que operan en mercados emergentes suelen experimentar un aumento de entre 0.25 y 0.50 puntos en su Beta debido a estos factores.

Como startup en etapa temprana, este proyecto enfrenta riesgos inherentes a la falta de una base de clientes establecida y de ingresos recurrentes probados. Esto justifica un ajuste hacia el rango superior del valor de Beta, posicionándolo en 1,5.

\paragraph{Prima de riesgo para Argentina}

Al obtener la prima de riesgo para Argentina, una referencia destacada es el trabajo del profesor Pablo Fernández de IESE Business School, quien publica regularmente estudios sobre las primas de riesgo de mercado a nivel mundial.

En su informe de 2023, Fernández estima la prima de riesgo de mercado para Argentina en torno al 7\%.

Este valor se obtiene considerando la diferencia entre el rendimiento promedio del mercado accionario argentino y la tasa libre de riesgo, ajustada por factores específicos del país.

\[Rm - Rf = 7\%\]

\paragraph{Riesgo País}

El riesgo país para Argentina a la fecha de elaboración de este trabajo se ubica en 578 puntos básicos.

\paragraph{Cálculo del costo del capital}

Al no tener fuentes de financiamiento y hacer frente solo con aportes de capital la tasa de descuento es igual al costo del capital.

\[Ke = 4,6\% + 1,5 \times 7\% + 5,78\% = 20,88\%\]

\subsubsection{Evaluación monetaria del proyecto}

\subparagraph{Duración}

La industria de la inteligencia artificial se encuentra en un momento de crecimiento explosivo, con cambios de paradigma que requieren repensar modelos de negocio constantemente. De esta manera se ha adjudicado una vida útil de 7 años al desarrollo del software de hiperpersonalización. Se tiene en cuenta una actualización durante el año 3 correspondiente al 50\% de la inversión inicial.


\clearpage


\paragraph{Estados de resultados proforma}

\begin{center}
\captionof{table}{Estados de resultados proforma}
\footnotesize
\begin{adjustbox}{width=\textwidth,center}
\begin{tabular}{l|rrrrrrrr}
\hline
\textbf{Año} & \textbf{0} & \textbf{1} & \textbf{2} & \textbf{3} & \textbf{4} & \textbf{5} & \textbf{6} & \textbf{7} \\
\hline
\multicolumn{9}{l}{\textbf{Ventas}} \\
\hline
Ventas netas & & \$ 220,000 & \$ 362,000 & \$ 582,000 & \$ 542,000 & \$ 602,000 & \$ 684,000 & \$ 756,000 \\
\hline
\multicolumn{9}{l}{\textbf{Costo de ventas}} \\
\hline
Costo de ventas & \$- & \$ 68,000 & \$ 98,800 & \$ 152,400 & \$ 118,371 & \$ 133,145 & \$ 150,538 & \$ 169,145 \\
\hline
Margen bruto & \$- & \$ 152,000 & \$ 263,200 & \$ 429,600 & \$ 423,629 & \$ 468,855 & \$ 533,462 & \$ 586,855 \\
\hline
\multicolumn{9}{l}{\textbf{Gastos operativos}} \\
\hline
Ventas, generales y administrativas & \$ 2,000 & \$ 212,456 & \$ 268,468 & \$ 314,364 & \$ 357,316 & \$ 360,316 & \$ 364,416 & \$ 391,032 \\
\hline
Gastos operativos totales & \$ 2,000 & \$ 212,456 & \$ 268,468 & \$ 314,364 & \$ 357,316 & \$ 360,316 & \$ 364,416 & \$ 391,032 \\
\hline
Resultado operativo & -\$ 2,000 & -\$ 60,456 & -\$ 5,268 & \$ 115,236 & \$ 66,313 & \$ 108,539 & \$ 169,046 & \$ 195,823 \\
\hline
\multicolumn{9}{l}{\textbf{Amortizaciones}} \\
\hline
Computadoras & \$- & \$ 6,000 & \$ 8,000 & \$ 9,000 & \$ 5,000 & \$ 3,000 & \$ 2,000 & \$ 1,000 \\
\hline
\multicolumn{9}{l}{\textbf{Otros ingresos y gastos}} \\
\hline
Renta extraordinaria & & \$- & \$ 3,000 & \$ 1,000 & \$ 500 & \$ 1,000 & \$ 500 & \\
\hline
Resultado antes de IG & -\$ 2,000 & -\$ 66,456 & -\$ 13,268 & \$ 109,236 & \$ 62,313 & \$ 106,039 & \$ 168,046 & \$ 195,323 \\
\hline
IG & -\$ 700 & -\$ 23,260 & -\$ 4,644 & \$ 38,233 & \$ 21,810 & \$ 37,114 & \$ 58,816 & \$ 68,363 \\
\hline
Resultado neto & -\$ 2,000 & -\$ 66,456 & -\$ 13,268 & \$ 71,003 & \$ 40,504 & \$ 68,926 & \$ 109,230 & \$ 126,960 \\
\hline
\end{tabular}
\end{adjustbox}
\end{center}

Hasta el año 2 el negocio no es sustentable y necesita del aporte de capital para sostener los costos. 


\clearpage


\paragraph{Balances proforma}

\begin{center}
\captionof{table}{Balances proforma}
\footnotesize
\begin{adjustbox}{width=\textwidth,center}
\begin{tabular}{l|rrrrrrrr}
\hline
\textbf{Año} & \textbf{0} & \textbf{1} & \textbf{2} & \textbf{3} & \textbf{4} & \textbf{5} & \textbf{6} & \textbf{7} \\
\hline
\multicolumn{9}{l}{\textbf{Activo corriente}} \\
\hline
Caja y bancos & \$- & \$ 20,000 & \$ 8,732 & \$ 85,735 & \$ 125,239 & \$ 197,165 & \$ 308,395 & \$ 433,355 \\
Provisión de pago de IVA & \$- & \$ 4,165 & \$ 6,997 & \$ 11,319 & \$ 10,765 & \$ 11,980 & \$ 13,602 & \$ 15,072 \\
Provisión de pago de IG & & & -\$ 4,644 & \$ 38,233 & \$ 21,810 & \$ 37,114 & \$ 58,816 & \$ 68,363 \\
\hline
Total activo corriente & \$- & \$ 24,165 & \$ 11,085 & \$ 135,287 & \$ 157,814 & \$ 246,259 & \$ 380,813 & \$ 516,790 \\
\hline
\multicolumn{9}{l}{\textbf{Activo fijo}} \\
\hline
Computadoras & \$- & \$ 12,000 & \$ 10,000 & \$ 4,000 & \$ 5,000 & \$ 2,000 & \$- & \$ 2,000 \\
\hline
Total activo fijo & \$- & \$ 12,000 & \$ 10,000 & \$ 4,000 & \$ 5,000 & \$ 2,000 & \$- & \$ 2,000 \\
\hline
\multicolumn{9}{l}{\textbf{Activo intangible}} \\
\hline
Código propietario & \$ 39,800 & \$ 39,800 & \$ 39,800 & \$ 54,800 & \$ 54,800 & \$ 54,800 & \$ 54,800 & \$ 54,800 \\
\hline
\textbf{Total Activo} & \textbf{\$ 39,800} & \textbf{\$ 75,965} & \textbf{\$ 60,885} & \textbf{\$ 194,087} & \textbf{\$ 217,614} & \textbf{\$ 303,059} & \textbf{\$ 435,613} & \textbf{\$ 573,590} \\
\hline
\hline
\multicolumn{9}{l}{\textbf{Pasivos corrientes}} \\
\hline
IVA a pagar & \$- & \$ 4,165 & \$ 6,997 & \$ 11,319 & \$ 10,765 & \$ 11,980 & \$ 13,602 & \$ 15,072 \\
IG & & & -\$ 4,644 & \$ 38,233 & \$ 21,810 & \$ 37,114 & \$ 58,816 & \$ 68,363 \\
\hline
Total Pasivo Corriente & \$- & \$ 4,165 & \$ 2,353 & \$ 49,552 & \$ 32,575 & \$ 49,094 & \$ 72,418 & \$ 83,435 \\
\hline
\hline
\multicolumn{9}{l}{\textbf{Patrimonio}} \\
\hline
Capital social & \$ 41,800 & \$ 140,256 & \$ 140,256 & \$ 155,256 & \$ 155,256 & \$ 155,256 & \$ 155,256 & \$ 155,256 \\
Utilidades retenidas & & -\$ 2,000 & -\$ 68,456 & -\$ 81,724 & -\$ 10,721 & \$ 29,783 & \$ 98,709 & \$ 207,939 \\
Resultado del período & -\$ 2,000 & -\$ 66,456 & -\$ 13,268 & \$ 71,003 & \$ 40,504 & \$ 68,926 & \$ 109,230 & \$ 126,960 \\
\hline
Total Patrimonio & \$ 39,800 & \$ 71,800 & \$ 58,532 & \$ 144,535 & \$ 185,039 & \$ 253,965 & \$ 363,195 & \$ 490,155 \\
\hline
\hline
\textbf{Total P + PN} & \textbf{\$ 39,800} & \textbf{\$ 75,965} & \textbf{\$ 60,885} & \textbf{\$ 194,087} & \textbf{\$ 217,614} & \textbf{\$ 303,059} & \textbf{\$ 435,613} & \textbf{\$ 573,590} \\
\hline
\end{tabular}
\end{adjustbox}
\end{center}

\clearpage


\paragraph{Flujo de fondos}

\begin{center}
\captionof{table}{Flujo de fondos}
\footnotesize
\begin{adjustbox}{width=\textwidth,center}
\begin{tabular}{l|rrrrrrrr}
\hline
\textbf{Año} & \textbf{0} & \textbf{1} & \textbf{2} & \textbf{3} & \textbf{4} & \textbf{5} & \textbf{6} & \textbf{7} \\
\hline
\multicolumn{9}{l}{\textbf{Caja de operaciones}} \\
\hline
Ingreso neto & -\$ 2,000 & -\$ 66,456 & -\$ 13,268 & \$ 71,003 & \$ 40,504 & \$ 68,926 & \$ 109,230 & \$ 126,960 \\
\hline
\multicolumn{9}{l}{\textbf{Ajustado por:}} \\
\hline
Amortizaciones & \$- & \$ 6,000 & \$ 8,000 & \$ 9,000 & \$ 5,000 & \$ 3,000 & \$ 2,000 & \$ 1,000 \\
Cambios en cuentas & \$- & \$ 20,000 & \$ 8,732 & \$ 85,735 & \$ 125,239 & \$ 197,165 & \$ 308,395 & \$ 433,355 \\
\hline
\textbf{Caja de operaciones} & \textbf{-\$ 2,000} & \textbf{-\$ 40,456} & \textbf{\$ 3,464} & \textbf{\$ 165,738} & \textbf{\$ 170,743} & \textbf{\$ 269,091} & \textbf{\$ 419,625} & \textbf{\$ 561,315} \\
\hline
\hline
\multicolumn{9}{l}{\textbf{Caja de inversiones}} \\
\hline
Compraventa de bienes & \$- & -\$ 18,000 & -\$ 6,000 & -\$ 3,000 & -\$ 6,000 & \$- & \$- & -\$ 3,000 \\
Software & -\$ 39,800 & & & -\$ 15,000 & & & & \\
\hline
\textbf{Caja de inversiones} & \textbf{-\$ 39,800} & \textbf{-\$ 18,000} & \textbf{-\$ 6,000} & \textbf{-\$ 18,000} & \textbf{-\$ 6,000} & \textbf{\$-} & \textbf{\$-} & \textbf{-\$ 3,000} \\
\hline
\hline
\multicolumn{9}{l}{\textbf{Cambio neto en caja}} \\
\hline
& -\$ 41,800 & -\$ 58,456 & -\$ 2,536 & \$ 147,738 & \$ 164,743 & \$ 269,091 & \$ 419,625 & \$ 558,315 \\
\hline
\hline
\multicolumn{9}{l}{\textbf{Cambio en capital de trabajo}} \\
\hline
& \$- & \$ 20,000 & \$ 8,732 & \$ 85,735 & \$ 125,239 & \$ 197,165 & \$ 308,395 & \$ 433,355 \\
\hline
\textbf{CAPEX} & -\$ 39,800 & -\$ 18,000 & -\$ 6,000 & -\$ 18,000 & -\$ 6,000 & \$- & \$- & -\$ 3,000 \\
\hline
\end{tabular}
\end{adjustbox}
\end{center}




\paragraph{Valuación y viabilidad del proyecto}

Tomando una duración de 7 años, con una inversión inicial de USD 41,800 al inicio y un costo promedio del capital de 0,2088 anual procedemos al cálculo del valor actual neto de los flujos de fondo proyectados. 

\paragraph{Flujos de fondos proyectados}

Para el cálculo del VAN se necesitan los flujos de caja libre para la empresa conocidos como FCFF (free cash flow to the firm) que se computan como la suma de los resultados netos, las amortizaciones, las diferencias de capital de trabajo y los aportes de capital. 

\paragraph{Cálculo del VAN}

\[VAN = \sum_{t=0}^{n} \frac{FCF_t}{(1+r)^t} - \text{Inversión inicial}\]

\paragraph{Cálculo del TIR}

\[TIR = \text{Tasa que hace el VAN igual a cero}\]

\clearpage

\begin{center}
\captionof{table}{Flujo de caja libre para la empresa (FCFF)}
\footnotesize
\begin{adjustbox}{width=\textwidth,center}
\begin{tabular}{l|rrrrrrrr}
\hline
\textbf{Año} & \textbf{0} & \textbf{1} & \textbf{2} & \textbf{3} & \textbf{4} & \textbf{5} & \textbf{6} & \textbf{7} \\
\hline
\multicolumn{9}{l}{\textbf{Resultados}} \\
\hline
Resultado neto & -\$ 2,000 & -\$ 66,456 & -\$ 13,268 & \$ 71,003 & \$ 40,504 & \$ 68,926 & \$ 109,230 & \$ 126,960 \\
\hline
Amortizaciones & \$- & \$ 6,000 & \$ 8,000 & \$ 9,000 & \$ 5,000 & \$ 3,000 & \$ 2,000 & \$ 1,000 \\
\hline
Diferencia en Capital de Trabajo & \$- & \$ 20,000 & \$ 8,732 & \$ 85,735 & \$ 125,239 & \$ 197,165 & \$ 308,395 & \$ 433,355 \\
\hline
\textbf{CAPEX} & \textbf{-\$ 39,800} & \textbf{-\$ 18,000} & \textbf{-\$ 6,000} & \textbf{-\$ 18,000} & \textbf{-\$ 6,000} & \textbf{\$-} & \textbf{\$-} & \textbf{-\$ 3,000} \\
\hline
\textbf{FCFF} & \textbf{-\$ 41,800} & \textbf{-\$ 58,456} & \textbf{-\$ 2,536} & \textbf{\$ 147,738} & \textbf{\$ 164,743} & \textbf{\$ 269,090} & \textbf{\$ 419,625} & \textbf{\$ 558,315} \\
\hline
\end{tabular}
\end{adjustbox}
\end{center}

\paragraph{VAN}

Con las tasas calculadas y este flujo de fondos se arriba a un VAN USD 455,720.32 , por lo tanto al ser este número positivo se demuestra la viabilidad del proyecto. 

